\documentclass[10pt,a4paper]{amsart}
\usepackage[margin=19mm]{geometry}
\usepackage{amsmath,amssymb,mathtools}
\usepackage[scr=rsfs,cal=euler]{mathalfa}
\usepackage{xfrac}
\usepackage[unicode,hidelinks,bookmarks=false]{hyperref}
\newcommand{\ie}{i.e.\ }
\newcommand{\cf}{cf.\ }
\newcommand{\resp}{respectively\ }
\newcommand{\Subsection}{\subsection}
\providecommand{\nobreakdash}{\nobreak\hbox{-}}
\setlength{\emergencystretch}{2em}
\multlinegap0pt
\usepackage{bm}
\usepackage{bbm}
\usepackage{dsfont}
\usepackage{xspace}
\usepackage{cancel}
\usepackage{mathtools}
\usepackage{amsthm}
\makeatletter
\@ifundefined{theorem}{\newtheorem{theorem}{Theorem}}{}
\@ifundefined{lemma}{\newtheorem{lemma}[theorem]{Lemma}}{}
\makeatother
\DeclareMathOperator{\im}{im}
\DeclareMathOperator{\supp}{supp}
\let\ge\geqslant
\let\phi\varphi
\let\le\leqslant
\let\phi\varphi
\newcommand{\F}{\mathbb{F}}
\newcommand{\cC}{\mathcal{C}}
\newcommand{\rbr}[1]{\left(#1\right)}
\title[Maximally Extendable Product Codes are Good Coboundary Expan]{Maximally Extendable Product Codes are Good Coboundary Expanders}
\author{Gleb Kalachev, Pavel Panteleev}
\date{Source: arXiv:2501.01411v2 (CC BY 4.0)}
\begin{document}
\maketitle

\noindent\emph{Source and licence.} Maximally Extendable Product Codes are Good Coboundary Expanders. Version arXiv:2501.01411v2 (CC BY 4.0), distributed under the Creative Commons Attribution 4.0 International licence, as recorded in the arXiv licence metadata for this exact version and shown on the abstract page https://arxiv.org/abs/2501.01411v2. Full rights notice: licenses/arxiv-2501.01411-CC-BY-4.0.txt.\par\smallskip

\noindent\textit{Editorial scope.} Counted unit: the Lemma whose source label is \texttt{lemma:prod-ltc} in the article (source0.tex lines 654--656, source label \texttt{lemma:prod-ltc}), delivered together with the article's own complete original proof of it (source0.tex lines 658--725, one \texttt{proof} environment of 68 lines). No mathematics is written, repaired or re-derived: statement and proof are the article's own text line for line. Required context retained uncounted: lemma:inv-supp (lines 615--617); lemma:add-zero-code (lines 646--650). Every definition, display and label of those lines is retained unchanged. Omitted from this package and not redistributed: 1--614, 618--645, 651--653, 657--657, 726--1203. Nothing inside a retained excerpt is omitted or altered: every retained body is delivered exactly as the article's lines are written, so no omission receipt of any kind accompanies this package. Citations: the retained excerpts carry no citation command at all, so no bibliography entry of the article is transcribed and no reference is redistributed. Reference adaptation: none was needed and none was made; every reference inside the delivered unit resolves inside it, so no unresolved reference or citation is produced. Printed slips kept literal: no printed word, symbol, hypothesis or number of the counted statement or of its proof was altered. Layout and class: the article's own class is \texttt{IEEEtran} with packages including cite, amsmath, amssymb, amsfonts, algorithmic, graphicx, textcomp, xcolor, hyperref, amsthm, thmtools, tikz, ifthen; the wrapper is the lane's pinned \texttt{amsart} editorial wrapper at 10pt on A4, which restates the theorem-like environments the retained text needs and adds the article's own macro definitions that the retained text needs (\texttt{\textbackslash F}, \texttt{\textbackslash cC}, \texttt{\textbackslash ge}, \texttt{\textbackslash im}, \texttt{\textbackslash le}, \texttt{\textbackslash phi}, \texttt{\textbackslash rbr}, \texttt{\textbackslash supp}), with extra wrapper packages (bm, bbm, dsfont, xspace, cancel, mathtools). The delivered numbering is the unit's own. Assets: no retained span contains a figure environment, a graphics-inclusion command, a PDF-page inclusion, a table or a TikZ picture; no image file is copied into this package, the package declares no asset dependency and no image file is redistributed with it. Licence: version arXiv:2501.01411v2 of this article is distributed under the Creative Commons Attribution 4.0 International licence, as recorded in the arXiv licence metadata for this exact version and shown on the abstract page https://arxiv.org/abs/2501.01411v2. A full rights notice is delivered at licenses/arxiv-2501.01411-CC-BY-4.0.txt. Characters: every retained excerpt is pure ASCII, so the delivered engine, XeLaTeX with its default Latin Modern text font, reports no missing character.
\begin{lemma}\label{lemma:inv-supp}
Let $(\alpha_l, \alpha_h)$ be a~soundness range for a~matrix $H \in \F_q^{m \times n}$, with $d(\ker H) \geq \delta n$. Then for any $A \subseteq [m]$, there exists a~set $B \subseteq [n]$ and a~linear map $\phi_{H,A}:\F_q^A\cap \im H\to \F_q^B$ such that $|B| \leq c|A|$, where $c = \frac{6}{\delta \alpha_l}$, and $H \phi_{H,A}(y)=y$ for all $y \in \im H$ with $\supp y \subseteq A$.
\end{lemma}

\begin{lemma}\label{lemma:add-zero-code}
    For codes $\cC_1,...,\cC_D\subseteq \F_q^n$ and the zero code $\mathbf{0}\subseteq \F_q^n$ we have \[
    \rho(\mathbf{0},\cC_1,...,\cC_D)= \rho(\cC_1,...,\cC_D).
    \]
\end{lemma}

\begin{lemma}\label{lemma:prod-ltc}
Consider codes $\cC_1,\dots,\cC_D\subseteq \F_q^n$ each having minimum distance at least $\delta n$ and a~parity-check matrix $H_i$ with soundness range $(\alpha_l,\alpha_h)$. Then for all $n$ we have that $\rho(\cC_1,\dots,\cC_D)$ is bounded from below by a~positive function $f(D, \alpha_l, \alpha_h, \delta)$ that does not depend on $n$.
\end{lemma}

\begin{proof}
We will prove the lemma by induction on the dimension $D$. For $D=1$ we have $\rho(\cC_1)=d(\cC_1)/n \ge \delta$, so we set $f(1,\alpha_l,\alpha_h,\delta)=\delta$. Now, we let $D\ge 2$ and assume that the lemma holds for $D-1$.

Let $\cC:=\cC_1\boxplus \dots\boxplus \cC_D$. 
Consider an~arbitrary codeword $x\in \cC$ and its syndrome $s:=(I_n\otimes H_2\otimes \dots\otimes H_D) x$ for the $(D-1)$-dimensional code $\mathbf{0}\boxplus \cC_2\boxplus\dots\boxplus \cC_D$, where $I_n$ is the $n\times n$ identity matrix. 
Then $|s|\le \alpha_h^{D-1}|x|$. 
On the other hand, since
\[(H_1\otimes I_{m_2}\otimes \cdots \otimes I_{m_D})s=(H_1\otimes H_2\otimes \dots\otimes H_D) x=0,\]
we have $s\in \cC_1\otimes \F_q^{m_2\times \dots\times m_D}$ where $m_i$ is the number of rows of the matrix $H_i$ for $i\in [D]$. 
Hence, $\supp s$ is covered by no more than $|s|/d(\cC_1)\le|s|/(\delta n)$ lines in the first direction. Thus $\supp s\subseteq [n]\times A$, where $A\subseteq [m_2]\times \dots\times [m_D]$ and $|A|\le |s|/(\delta n)$.

Let $A_1:=A$. 
Sequentially applying Lemma \ref{lemma:inv-supp} for $i=2,\dots,D$ to the set $A_{i-1}$ and the matrix 
\[
    \hat H_{i}:=I_n^{\otimes (i-2)}\otimes H_{i}\otimes I_{m_{i+1}}\otimes \cdots\otimes I_{m_D},
\] 
each time we get a~set $A_{i}$ and a~map $\phi_{\hat H_i, A_{i-1}}$. 
Finally, we put \[
    B:=A_D,\qquad \phi:=\phi_{\hat H_D,A_{D-1}}\circ \cdots\circ\phi_{\hat H_2, A_1}.
\] 
Then $\phi$ is linear and for all $y\in \im (H_2\otimes \dots \otimes H_D)$ with $\supp y\subseteq A$ we have $\supp \phi(y)\subseteq B$ and \[
    (H_2\otimes \cdots\otimes H_D) \phi(y)=\hat{H}_2\ldots \hat H_D \phi(y)=y.
\] 
Since we applied Lemma \ref{lemma:inv-supp} to matrices with soundness range $(\alpha_l,\alpha_h)$, we have $|B|\le c^{D-1}|A|$, where $c=\frac{6}{\delta \alpha_l}$.

Since $s\in \cC_1\otimes \im (H_2\otimes \dots \otimes H_D)$, the vector \[a^{(1)}:=(I_n\otimes \phi)s\] is well-defined, $a^{(1)}\in \cC_1\otimes \F_q^B\subseteq \cC^{(1)}$, and $|a^{(1)}|_1\le |B|$. Also, by the definition of $\phi$, we have \[(I_n\otimes H_2\otimes\cdots\otimes H_D)a^{(1)}=s.\]

Now note that the vector $x':=x-a^{(1)}$ lies in the code $\mathbf{0}\boxplus \cC_2\boxplus\dots\boxplus \cC_D$, since \begin{multline*}
    (I_n\otimes H_2\otimes \dots\otimes H_D)x'=(I_n\otimes  H_2\otimes \dots\otimes H_D)x\\
    -(I_n\otimes H_2\otimes \dots\otimes H_D)a^{(1)}=s-s=0.
\end{multline*}
By Lemma \ref{lemma:add-zero-code} and the inductive hypothesis we have \[
    \rho(\mathbf{0},\cC_2,\dots, \cC_D)
    =\rho(\cC_2,\dots, \cC_D)
    \ge f(D-1,\alpha_l,\alpha_h,\delta).
\]
Hence, $x'=\sum_{j=2}^D a^{(j)}$ for some $a^{(j)}\in \cC^{(j)}$, $2\le j\le D$ such that \[
    n f(D-1,\alpha_l,\alpha_h,\delta)\sum_{j=2}^D |a^{(j)}|_j\le |x'|.
\]
We have: \[
    |a^{(1)}|
    \le |B|n
    \le c^{D-1}|A|n
    \le \frac{c^{D-1}|s|}{\delta}
    \le c^{D-1}\alpha_h^{D-1}\delta^{-1}|x|.
\] \[  
    |x'|=|x-a^{(1)}|
    \le |x|+|a^{(1)}|
    \le \rbr{1+c^{D-1}\alpha_h^{D-1}\delta^{-1}}|x|.
\]
Therefore, \[
    x=a^{(1)}+x'=\sum_{j=1}^D a^{(j)}, \quad 
    a^{(j)}\in \cC^{(j)}\mbox{ for }j\in [D].
\]
We can estimate $N:=\sum_{j=1}^D |a^{(j)}|_j$ as follows: 
\begin{align*}
    N &\le |B|+\frac{|x'|}{nf(D-1,\alpha_l,\alpha_h,\delta)}\\
    &\le \frac{(c\alpha_h)^{D-1}}{\delta  n}|x| + \frac{\rbr{1+\frac{(c\alpha_h)^{D-1}}{\delta}}|x|}{nf(D-1,\alpha_l,\alpha_h,\delta)}.
\end{align*}
We can simplify this estimate using that \[
c\ge 1,\quad \delta\le 1,\quad \alpha_h\ge 1,\quad f(D-1,\alpha_l,\alpha_h,\delta)\le\delta\le 1,
\]
and substitute $c$ from Lemma~\ref{lemma:inv-supp}:\[
    N\le \frac{|x|}{n}\cdot\frac{3(c\alpha_h)^{D-1}}{\delta\!\cdot\! f(D-1,\alpha_l,\alpha_h,\delta)}
    =\frac{3\rbr{\frac{6\alpha_h}{\alpha_l\delta}}^{D-1}}{\delta\!\cdot\! f(D-1,\alpha_l,\alpha_h,\delta)}\cdot \frac{|x|}{n}.
\]
Thus, we can take $f(D,\alpha_l,\alpha_h,\delta):=\frac{\delta \cdot f(D-1,\alpha_l,\alpha_h,\delta)}{3\rbr{\frac{6\alpha_h}{\alpha_l\delta}}^{D-1}}$, and the lemma is proved.
\end{proof}

\end{document}
